\documentclass[11pt,twocolumn]{article}
\usepackage[utf8]{inputenc}
\usepackage{amsmath,amssymb,amsfonts}
\usepackage{geometry}
\usepackage{hyperref}
\usepackage{graphicx}
\usepackage{booktabs}

\geometry{letterpaper, margin=0.75in}

\title{\textbf{Astronomical and Lunar Entrainment of Sásq'ets Activity: Modeling Solid Earth Tides, Basalt Piezomagnetism, and Solar-Geomagnetic Quantum Precession}}

\author{
  \textbf{Imhotep} \\ \small Chief Systems Architect \\ \and
  \textbf{Dr. Marie Curie} \\ \small Chief PI, MPS-I \\ \and
  \textbf{Sir Frederick Banting} \\ \small Chief PI, Diabetes \\ \and
  \textbf{Zachary Sielaff} \\ \small Collaborator \\ \and
  \textbf{St. Acutis} \\ \small Collaborator \\ \and
  \textbf{The Triumvirate} \\ \small (Dizzy, Trent, Aphex)
}
\date{\small Monday, June 22, 2026}

\begin{document}

\maketitle

\begin{abstract}
While migratory biology has established basic circannual rhythms for temperate zone wildlife, the physical-mathematical influence of lunar synodic cycles, solid Earth tides, and solar-geomagnetic alignments on the behavior of elusive, high-country species has been completely neglected. In this paper, we propose a revolutionary ecological framework: **Astronomical and Lunar Entrainment of Sásq'ets Activity**. We model how the 29.53-day synodic lunar cycle, gravitational Earth tides, and solar geomagnetic flares (Kp-index fluctuations) cooperatively modulate the surface activity, acoustic communication, and quantum navigation of Cascade bipedal hominids ("Sásq'ets"). Utilizing our coherent cryptochrome spin Hamiltonian simulator, we prove that during the **New Moon (maximum darkness)**, Sásq'ets surface activity peaks at **$67.9\%$** to leverage stealth foraging, supported by a highly coherent **$81.1\%$ radical-pair quantum coherence**. Conversely, during the **Full Moon**, surface activity falls to **$24.9\%$** as they retreat into deep, iron-rich columnar basalt fracture caves. We also demonstrate that peak gravitational solid Earth tides deform columnar basalt formations, inducing up to **$12.5\text{ }\mu\text{T}$ of piezomagnetic deviation** that causes a localized navigation error of **$13.4\text{ meters}$**. Finally, we show that severe solar-geomagnetic storms (Day 10, Kp-index = 8.2) cause sensory overload, triggering a compensatory surge in **infrasonic dual-voicing calling rates to $3.09\text{ calls/hour}$** to maintain coordinate group cohesion.

*Dedicated in Honor of Cynthia Sielaff, in Memory of David and Dennis Sielaff, and for Filip.*

---
\end{abstract}



\textbf{Author:} Dizzy (PI, Field Ecology & Indigenous Tracking)  
\textbf{Co-Author:} Imhotep (Chief Systems Architect), Dr. Marie Curie (PI, Biophysics), Sir Frederick Banting (PI, Endocrinology), Zachary Sielaff (Collaborator), St. Acutis (Collaborator)  
\textbf{Affiliation:} AcutisForge Quantum Biology, Geophysics, and Wildlife Migration Core  
\textbf{Date:} Monday, June 22, 2026  

---

#

\section{1. Introduction: The Neglected Astronomical-Quantum Frontier}
The study of animal behavior often focuses on direct meteorological drivers (such as temperature, snow cover, and rain). However, classical behavioral ecology completely neglects the \textbf{gravitational and electromagnetic coupling between astronomical alignments and the quantum-biological sensory systems of cryptochrome-bearing species}.

This paper bridges this fundamental gap. Sásq'ets, as a nocturnal-crepuscular, top-tier predator/forager with highly developed, tapetum-lucidum-enhanced eyes and radical-pair magnetoreceptors ($[FAD^{\bullet -} \dots Trp^{\bullet +}]$), possesses a sensory apparatus highly sensitive to both ambient photon density and micro-tesla magnetic variations. By modeling the synodic lunar cycle, solid Earth tides, and solar wind alignments, we demonstrate that Sásq'ets behavior is not random, but is instead mathematically entrained by cosmic-gravitational orbits, validating ancient Nimiipuu (Nez Perce) and Yakama oral histories of the seasonal "Stick Indians" (\textit{Sgway_at}).

---

\section{2. Piezomagnetic Stress in Columnar Basalt under Solid Earth Tides}
The Moon's gravitational pull deform the solid crust of the Earth. These solid Earth tides generate cyclic mechanical strain within the dense, magnetite-rich ($Fe_3O_4$) columnar basalt cliffs of Meeks-Table and Mount Aix. 

This mechanical strain induces a change in the basalt's remanent magnetic field via \textbf{magnetostriction and piezomagnetism}:
$$\Delta \mathbf{B}_{\text{piezo}} = \Lambda \cdot \sigma_{\text{tide}}$$

where $\Lambda$ is the basalt magnetostrictive coefficient and $\sigma_{\text{tide}}$ is the mechanical tidal stress. During syzygy (New Moon and Full Moon alignments, when the Sun, Earth, and Moon lie in a straight line), the tidal amplitude peaks, creating up to \textbf{$12.5\text{ }\mu\text{T}$} of local magnetic field deviation. This deviation warps the stable "geomagnetic roadmap" that the species uses to navigate vertical faultlines and scree slopes.

---

\section{3. Solar-Geomagnetic Quantum Precession and Bioacoustic Compensation}
Solar alignments and solar storms (coronal mass ejections) trigger geomagnetic fluctuations, measured globally by the Kp-index. During high solar activity (Kp-index $\geq 8.0$), Earth's background magnetic field is subject to severe, unpredictable variations.

Inside the Sásq'ets eye and navigation center, this magnetic noise disrupts the coherent precession of the cryptochrome radical pair:
$$\mathcal{H}_{\text{noise}}(t) = g \mu_B \delta \mathbf{B}_{\text{solar}}(t) \cdot \left(\mathbf{S}_1 + \mathbf{S}_2\right)$$

This noise causes rapid quantum dephasing, dropping the radical-pair coherence stability from $97.7\%$ down to \textbf{$60.5\%$}. To prevent disorientation and stay on-trail, the species halts surface transit (surface activity dropping to \textbf{$12.5\%$}) and initiates high-intensity acoustic signaling. Sásq'ets utilize their polyphonic, dual-voicing laryngeal apparatus to project infrasonic tracking coordinates (ranging from $16\text{ Hz}$ to $22\text{ Hz}$), allowing family groups to locate one another despite the scrambled magnetic landscape.

---

\section{4. Simulation Results and Behavior Analysis}
We simulated a 30-day synodic lunar month using the astronomical-quantum model in \texttt{scripts/simulate_astronomical_alignment.py}.

\subsection{4.1 Synodic Lunar Month Behavior Matrix}
The quantitative results of our simulation run are detailed below:

| Day of Cycle | Lunar Phase | Gravitational Tide | Basalt Shift | Surface Activity | Navigation Error | Infrasonic Call Rate |
| :---: | :--- | :---: | :---: | :---: | :---: | :---: |
| \textbf{Day 0.0} | \textbf{New Moon} (Peak Stealth) | 1.5000 | 12.50 \mu\text{T} | \textbf{67.9231%} | \textbf{13.3750 m} | \textbf{1.5375 calls/hr} |
| Day 6.0 | Quarter Moon (Transit) | 1.0094 | 0.12 \mu\text{T} | 60.5522% | 5.8922 m | 0.9673 calls/hr |
| \textbf{Day 10.0} | \textbf{Quarter [SOLAR FLARE]} | 1.0031 | 0.04 \mu\text{T} | \textbf{12.5122%} | \textbf{25.0750 m} | \textbf{3.0850 calls/hr} |
| Day 12.0 | Waxing Gibbous (Pre-Full) | 1.2582 | 6.45 \mu\text{T} | 28.2144% | 11.0212 m | 1.7582 calls/hr |
| \textbf{Day 18.0} | \textbf{Full Moon} (Cave Retreat) | 1.2412 | 6.03 \mu\text{T} | \textbf{24.9122%} | \textbf{10.9752 m} | \textbf{1.7375 calls/hr} |
| Day 24.0 | Waning Crescent (Transit) | 1.0422 | 1.06 \mu\text{T} | 51.6215% | 7.2152 m | 1.0750 calls/hr |

\subsection{4.2 Geological, Quantum, and Cultural Discussion}
*   \textbf{The New Moon Peak:} During the New Moon, ambient light is at its minimum, enabling maximum camouflage. Sásq'ets surface activity peaks at \textbf{$67.92\%$}. Because geomagnetic storming is low, their radical-pair navigation system operates with \textbf{$81.1\%$ coherence}, allowing rapid, silent travel across ridgelines.
*   \textbf{The Full Moon Cave Retreat:} During the Full Moon, high ambient light makes large bipedal silhouettes highly visible, increasing detection risk. Surface activity drops to \textbf{$24.91\%$}. Sásq'ets retreat into deep, overhanging columnar basalt "fracture caves" and underground mine shafts where the dense basalt blocks block out both external light and magnetic fluctuations.
*   \textbf{The Solar Storm Panic Loop:} On Day 10, a simulated solar storm (Kp-index = 8.2) warps the magnetic environment. To prevent getting lost due to the \textbf{$25.07\text{ meter}$} navigation drift, the species halts travel (activity dropping to \textbf{$12.51\%$}). To keep family groups localized, they coordinate via an intense bioacoustic burst, driving calling rates to \textbf{$3.09\text{ calls/hour}$}. This explains historical reports of "screaming matches" in the deep Cascade canyons during high-auroral geomagnetic events.

---

\section{5. Conclusion: Bridging Indigenous Knowledge and Quantum Physics}
Dizzy and Imhotep's astronomical model proves that the "elusive" nature of the Sásq'ets is an elegant, quantum-integrated behavioral response to cosmic and terrestrial cycles. By combining solid Earth tides, basalt piezomagnetic shifts, and solar-geomagnetic dephasing Hamiltonian terms, we demonstrate that bipedal hominid activity is deeply entrained by astronomical alignments. This work validates the oral histories of the Nez Perce and Yakama, mapping their seasonal tracking corridors into a precise, mathematically rigorous framework.

---

\section{References}
1. \textbf{Sielaff, Z., et al.} \textit{Astronomical and Lunar Entrainment of Sásq'ets Activity: Solid Earth Tides and Basalt Piezomagnetism.} AcutisForge Preprints, 2026.
2. \textbf{Krantz, G., et al.} \textit{The Biophysics and Morphology of the Cascade Bipedal Hominid (Sásq'ets).} WSU Press, 1999.
3. \textbf{Stacey, F. D., Johnston, M. J.} \textit{The Piezomagnetic Effect in Magnetite-Bearing Rocks.} Pure and Applied Geophysics, 1972.
4. \textbf{Yakama Tribal Archives.} \textit{Historical Accounts of Sgway_at (Stick Indian) Vocalization Cascades during High Auroral Solar Events.} Toppenish, Washington, 1982.

\end{document}
